\subsection{C. Koszul regularity: what survives, and what fails}

All rings in this subsection are commutative with identity.  For a finite
sequence $f=(f_1,\ldots,f_r)$ in $R$, let $K_R(f)$ be the differential
graded exterior algebra on $R^r$ in homological degree one, with
$d(e_i)=f_i$.  For an arbitrary $R$-module $M$, we call $f$
$M$-Koszul-regular if $H_q(K_R(f)\otimes_R M)=0$ for every $q>0$.
No finiteness is imposed on $M$.  We reserve ordinary $M$-regularity for
successive injectivity on $M/(f_1,\ldots,f_{i-1})M$, together with
$M/(f)M\ne0$.

\begin{proposition}[Proved: the literal degree-zero reading is false]\label{SA23-C-1}
The useful condition on cote 23, scan 053, is vanishing of Koszul
\emph{homology in positive degrees}.  If its transcription ``nul en
degr\'es $0$'' means vanishing of $H_0$, it says instead $(f)=R$.
For a proper ideal it is never satisfied.  If it means the degree-zero
term of the complex is zero, it forces $R=0$.
\end{proposition}
\begin{proof}
The degree-zero term is $R$, and the image of its incoming differential
is $(f)$; consequently $H_0(K_R(f))=R/(f)$.  These identities prove all
claims.  We do not silently emend a doubtful reading into a theorem.
\end{proof}

\begin{theorem}[Proved: Koszul regularity implies quasi-regularity]\label{SA23-C-2}
Put $I=(f_1,\ldots,f_r)$.  If $f$ is $M$-Koszul-regular, the natural
map of graded modules
\[
 (M/IM)[T_1,\ldots,T_r]\longrightarrow
 \bigoplus_{n\ge0}I^nM/I^{n+1}M,\qquad T_i\longmapsto f_i,
\]
is an isomorphism.  In particular, if $M=R$, $I/I^2$ is free of rank
$r$ over $R/I$, with the displayed generators as basis, and
$\operatorname{Sym}_{R/I}(I/I^2)\to\operatorname{gr}_I R$ is an
isomorphism.  Ordinary regular sequences are Koszul-regular.
\end{theorem}
\begin{proof}
Set $S=\mathbf Z[T_1,\ldots,T_r]$ and $J=(T_1,\ldots,T_r)$, and
make $M$ an $S$-module by $T_i m=f_i m$.  The Koszul resolution of
$\mathbf Z=S/J$ shows that
$\operatorname{Tor}^S_q(S/J,M)=0$ for $q>0$.  Each $J^n/J^{n+1}$ is
a finite free $S/J$-module, with its monomial basis; hence it too has
zero positive Tor against $M$.  The sequences
\[
 0\to J^n/J^{n+1}\to S/J^{n+1}\to S/J^n\to0
\]
show inductively that $\operatorname{Tor}^S_q(S/J^n,M)=0$ for
$q>0$.  Tensoring $0\to J^n\to S\to S/J^n\to0$ therefore identifies
$J^n\otimes_S M$ with its image $I^nM$ in $M$.  Tensoring
$0\to J^{n+1}\to J^n\to J^n/J^{n+1}\to0$ now gives
\[
 I^nM/I^{n+1}M\simeq(J^n/J^{n+1})\otimes_S M.
\]
The monomial bases give exactly the proposed graded map, including
its compatibility with multiplication when $M=R$.

For the last assertion use induction on the sequence length.  The
Koszul complex on $f_1,\ldots,f_i$ is the cone of multiplication by
$f_i$ on the preceding Koszul complex.  If the preceding complex has
only $H_0=M/(f_1,\ldots,f_{i-1})M$, the cone's positive homology
vanishes precisely when $f_i$ acts injectively there.
\end{proof}

\begin{theorem}[Proved: generators, flat descent, and composition]\label{SA23-C-3}
Suppose $I\subset R$ is generated by a Koszul-regular sequence of
length $r$.
\begin{enumerate}
\item Every other length-$r$ generating sequence of $I$ is
Koszul-regular.  In particular the property is invariant under
permutation.
\item For a flat map $R\to R'$, Koszul regularity ascends to $R'$.
For a faithfully flat map it also descends.  The same assertions hold
with arbitrary module coefficients $M$ and $M\otimes_R R'$.
\item The property that an ideal is locally generated by a finite
Koszul-regular sequence, and the corresponding property of an immersion
of schemes, descend for fpqc coverings.
\item If $f$ is Koszul-regular in $R$ and the images of
$g=(g_1,\ldots,g_s)$ are Koszul-regular in $R/I$, then $(f,g)$ is
Koszul-regular in $R$.  Thus compositions of Koszul-regular immersions
are Koszul-regular.
\end{enumerate}
\end{theorem}
\begin{proof}
For (1), localize at a prime $\mathfrak p$.  If $I\not\subset\mathfrak p$,
both Koszul complexes are contractible: choose $a_i$ with
$\sum a_i f_i=1$; exterior multiplication by $\sum a_i e_i$ is a
contracting homotopy, and likewise for the other generators.  If
$I\subset\mathfrak p$, Theorem~\ref{SA23-C-2} makes the images of
$f$ a basis of $I/I^2$ over $(R/I)_\mathfrak p$.  A second length-$r$
generating family is also a basis: its matrix is surjective on a
finite free module over a local ring and thus has unit determinant.
Writing $g=Uf$ in $R_\mathfrak p^r$, the determinant of $U$ is a
unit modulo $I$, hence a unit in $R_\mathfrak p$.  Exterior powers of
$U$ give an isomorphism of the Koszul complexes.  Vanishing after
localization at every prime proves vanishing over $R$.

For (2), there is an equality of complexes
\[
 (K_R(f)\otimes_R M)\otimes_R R'
 =K_{R'}(f\otimes1)\otimes_{R'}(M\otimes_R R').
\]
Flatness commutes with homology; faithful flatness detects zero modules.
For (3), finite type of the ideal and finite local freeness of its
conormal module descend by faithfully flat descent.  Locally downstairs,
lift a conormal basis.  Nakayama's lemma, applied to the finitely
generated ideal in the local ring, and then shrinking the affine
neighborhood, makes the lifts a generating family.  Upstairs these are
Koszul-regular by (1); (2) descends their acyclicity.  This argument is
local on the base and therefore applies to an fpqc cover.  The standard
faithfully flat descent facts used here are the descent of finite
presentation and of finite projective modules; no coherence or
Noetherian hypothesis is needed.

For (4), the augmentation $K_R(f)\to R/I$ is a quasi-isomorphism.
Tensor it with the bounded complex of finite free modules $K_R(g)$.
The result is again a quasi-isomorphism, from $K_R(f,g)$ to
$K_{R/I}(\overline g)$.  The latter is a resolution of $R/(f,g)$.
\end{proof}

\begin{proposition}[Proved: the valid converse transitivity]\label{SA23-C-4}
Let $R$ be local, let $I\subset J$ be finitely generated proper ideals,
and assume $I$ and $J$ are Koszul-regular.  If
\[
 0\longrightarrow I/IJ\longrightarrow J/J^2
 \longrightarrow J/(I+J^2)\longrightarrow0
\]
is split exact over $R/J$, then $J/I$ is Koszul-regular in $R/I$.
This holds locally for immersions without a Noetherian assumption.
\end{proposition}
\begin{proof}
Choose generators $f$ of $I$ lifting a basis of $I/I^2$; these are
Koszul-regular.  The split conormal sequence allows their images to be
extended to a basis of $J/J^2$, say by images of $g$.  Nakayama gives
$J=(f,g)$, and Theorem~\ref{SA23-C-3}(1) makes $(f,g)$
Koszul-regular.  The quasi-isomorphism
$K_R(f,g)\to K_{R/I}(\overline g)$ used in part (4) proves the
claim.  All generators and matrix inverses exist on a neighborhood of
the point, so the same proof applies to immersions.
\end{proof}

\begin{example}[Proved: a local nonpermutable sequence; refuted converse]\label{SA23-C-5}
Let $k$ be any field, put $D=k[x,y]_{(x,y)}$, and make $M=k(x)$ a
$D$-module by the homomorphism $D\to k(x)$ sending $y$ to zero.
This map is defined because the image at $y=0$ of a polynomial outside
$(x,y)$ is nonzero.  Form the idealization
\[
 R=D\ltimes M,\qquad
 (a,m)(b,n)=(ab,an+bm).
\]
We use $x,y$ also for $(x,0),(y,0)\in R$.  Then $R$ is local,
$(x,y)$ is an ordinary regular sequence, and $(y,x)$ is not.  For
$I=(y)\subset J=(x,y)$, both $J\subset R$ and $J/I\subset R/I$
are Koszul-regular and their conormal sequence is split exact; $I$ is
not Koszul-regular.  Thus the hoped-for implication
``the composite and the first immersion are Koszul-regular, hence the
second immersion is'' is false without additional hypotheses.
\end{example}
\begin{proof}
An element $(a,m)$ is a unit precisely when $a$ is a unit in $D$:
its inverse then is $(a^{-1},-a^{-2}m)$.  Thus the unique maximal ideal
is $(x,y)D\ltimes M$ and the residue field is $k$.
Multiplication by $x$ is injective on $D$ and bijective on $M$, so is
injective on $R$.  Since $xM=M$,
\[
 R/xR\simeq D/xD\simeq k[y]_{(y)}.
\]
The image of $y$ is a nonzerodivisor and the further quotient is $k$.
This proves regularity of $(x,y)$.  On the other hand, $y(0,1)=0$
and $(0,1)\ne0$, so $y$ is a zerodivisor.  In particular $(y,x)$
is not ordinary regular, although its Koszul complex is isomorphic
to that of $(x,y)$ and is acyclic in positive degrees.

The ideal $I=yD\ltimes0$ and its square $I^2=y^2D\ltimes0$ satisfy
$I/IJ\simeq k$, with basis the image of $y$.  The equality
$xM=M$ gives $J=(x,y)D\ltimes M$ and
$J^2=(x,y)^2D\ltimes M$, so $J/J^2\simeq k^2$ with basis $x,y$.
The conormal sequence is therefore
$0\to k\cdot y\to kx\oplus ky\to k\cdot x\to0$, visibly split.
Furthermore
$R/I\simeq k[x]_{(x)}\ltimes k(x)$ and multiplication by $x$ is
injective in this ring, with quotient $k$.  Thus $J/I$ is regular.
But $H_1(K_R(y))=(0:_R y)$ contains $0\ltimes k(x)$ and is nonzero.
No finite example could exhibit the claimed non-Noetherian phenomenon:
a finite ring is Noetherian.  This is an explicit infinite example,
not an assertion inferred from finite truncations.
\end{proof}

\begin{remark}[Known scope; exact obstruction]\label{SA23-C-6}
For finite modules over Noetherian local rings and sequences in the
maximal ideal, ordinary, Koszul, and quasi-regularity are equivalent
(Stacks Project, Lemma 15.31.7, tag \texttt{062D}; opened).
The word ``local'' cannot be dropped for sequences: in the Noetherian
ring $k[x,y]\times k[x,y]$ the sequence
$((x,1),(y,0))$ is regular, whereas the reversed sequence begins with
a zerodivisor.  Its final quotient is $k$.

The failure in Example~\ref{SA23-C-5} is the failure of Nakayama for an
infinitely generated homology module: $xM=M\ne0$.
For a sequence in the maximal ideal of a local ring $R$, if the
positive Koszul homology modules of every prefix are finite,
acyclicity of the full sequence does imply ordinary $R$-regularity.
Indeed, the cone long exact sequence shows
that multiplication by the last element is surjective on the positive
homology of the preceding complex; Nakayama kills it.  Repeating this
argument reduces to a one-element complex.  This gives a valid
finiteness replacement for the false unrestricted converse.  The
Noetherian case satisfies that replacement.  SGA 6 VII, Proposition
1.7(iii), explicitly retains a Noetherian hypothesis for precisely this
converse; Proposition 1.7(i)--(ii) records the two unrestricted
transitivities proved above.
\end{remark}

\subsubsection{C.3. Deligne's Sym--Tor characterization}

\begin{theorem}[Proved under the stated Noetherian hypothesis]\label{SA23-C-7}
Let $R$ be Noetherian and $I\subset R$ an ideal.  The following are
equivalent:
\begin{enumerate}
\item $I/I^2$ is projective over $R/I$ and
$\operatorname{Sym}_{R/I}(I/I^2)\to\operatorname{gr}_I R$ is an
isomorphism;
\item $I/I^2$ is projective over $R/I$ and the canonical graded
algebra map
$\bigwedge_{R/I}(I/I^2)\to\operatorname{Tor}^R_*(R/I,R/I)$ is an
isomorphism;
\item locally at every prime containing $I$, the ideal $I$ is generated
by a regular sequence.
\end{enumerate}
If $R$ is local, $I$ is proper, and $f_1,\ldots,f_r$ is a minimal
generating family, these are also equivalent to: the Koszul complex
$K_R(f)$ resolves $R/I$; and
$\operatorname{Ext}^q_R(R/I,R)=0$ for $q<r$.
These are the hypotheses of Deligne's proposition on cote 23, scan 063.
\end{theorem}
\begin{proof}
Every assertion in the first list is local along $V(I)$.  Indeed,
$I/I^2$ is finitely presented; being finite projective is equivalent
to being locally free, and kernels and cokernels of the displayed
maps vanish if and only if their localizations vanish.  We may thus
work in a local ring with $I$ proper.

Condition (1) makes a minimal generating family $f$ quasi-regular:
Nakayama identifies its images with a basis of the free conormal
module.  Quasi-regularity is equivalent here to regularity by the
Noetherian local criterion just cited.  Conversely regularity implies
(1) by Theorem~\ref{SA23-C-2}.  Regularity implies that $K_R(f)$ is a
resolution.  Tensoring it with $R/I$ sets every differential equal to
zero and gives condition (2), including its multiplicative structure.

For the remaining implication from (2), choose a degreewise finite
free differential graded resolution $P_\bullet\to R/I$ extending
the Koszul complex; the usual Tate construction supplies one by
successively adjoining generators to kill positive homology.
Noetherianity ensures that only finitely many generators are needed
in each degree.  The map
\[
 K_R(f)\otimes_R R/I\longrightarrow P_\bullet\otimes_R R/I
\]
induces exactly the canonical exterior-algebra-to-Tor map: in degree
one it is the identification of both sides with $I/I^2$, and it
respects products.  Under (2) its mapping cone becomes acyclic after
tensoring with $R/I$.  That cone is bounded below and degreewise
finite free.  Derived Nakayama makes it acyclic over $R$.  One can
verify this last step without an unbounded convergence assumption:
if it had nonzero homology, take its lowest nonzero degree $q$;
the lowest-degree edge of the hyper-Tor spectral sequence gives
$H_q(C)\otimes_R R/I$ in degree $q$ after derived tensoring.
This is nonzero by ordinary Nakayama, since $H_q(C)$ is finite and
$I$ lies in the maximal ideal.  This contradicts acyclicity.
Thus $K_R(f)\to P$ is a quasi-isomorphism and $K_R(f)$ is a resolution.
The Noetherian local criterion again makes $f$ regular.

Finally Lemma~\ref{SA23-C-8}, with $M=R$, shows that the indicated
Ext vanishing implies Koszul regularity and hence regularity here.
The converse follows by dualizing the self-dual length-$r$ Koszul
resolution.  The case $I=0$ has $r=0$ and all statements hold; for
$I=R$ only the first, local-along-$V(I)$ list applies vacuously.
\end{proof}

\begin{lemma}[Proved: the module form of the Ext implication]\label{SA23-C-8}
Let $R$ be any ring, let $I=(f_1,\ldots,f_r)$, and let $M$ be any
$R$-module.  If $\operatorname{Ext}^j_R(R/I,M)=0$ for $j<r$,
then $f$ is $M$-Koszul-regular.  If in addition $R$ is Noetherian
local, $I$ is proper, and $M$ is finite and nonzero, $f$ is an
ordinary $M$-regular sequence.  This proves the auxiliary lemma on
cote 23, scans 063--064, in its inherited Noetherian setting, and
gives a valid Koszul version without that setting.
\end{lemma}
\begin{proof}
First the Ext vanishing holds with $R/I$ replaced by any
$R/I$-module $N$.  For degree zero, embed
$\operatorname{Hom}_R(N,M)$ into
$\operatorname{Hom}_R(F,M)=0$, where $F\twoheadrightarrow N$ is
a free $R/I$-module.  Induct on the degree $j<r$, using
$0\to N_1\to F\to N\to0$.  Its Ext sequence places
$\operatorname{Ext}^j_R(N,M)$ between
$\operatorname{Ext}^{j-1}_R(N_1,M)=0$ and
$\operatorname{Ext}^j_R(F,M)=0$.  The latter vanishes because
Ext in the first argument takes direct sums to products; this
follows immediately by taking direct sums of projective
resolutions and using exactness of products of modules.

Every Koszul homology module $H_q(K_R(f))$ is annihilated by $I$:
exterior multiplication by $e_i$ is a homotopy between
multiplication by $f_i$ and zero.  Choose an injective resolution
$M\to E^\bullet$.  The double complex
$\operatorname{Hom}_R(K_q(f),E^p)$ has total degree $p+q$.
Its vertical-first spectral sequence computes the cohomology of
$\operatorname{Hom}_R(K_R(f),M)$, because each $K_q$ is free.
The horizontal-first spectral sequence has second page
\[
 E_2^{p,q}=
 \operatorname{Ext}^p_R(H_q(K_R(f)),M).
\]
It converges to the same total cohomology; there are only $r+1$
columns, so no unbounded-convergence issue arises.  Its terms with
$p<r$ vanish by the first paragraph.  Consequently the total
cohomology vanishes in degrees $j<r$.  Koszul self-duality identifies
that degree with $H_{r-j}(K_R(f)\otimes_R M)$, proving positive
Koszul acyclicity.  Under the additional Noetherian local hypotheses
the positive homology of each prefix is finite.  The cone and
Nakayama argument in Remark~\ref{SA23-C-6} gives successive
injectivity; ordinary Nakayama gives $M/IM\ne0$.
\end{proof}
